% GENERATED FILE. Edit canonical lessons, then run npm run generate:books.
% canonical-source-hash: fnv1a64:5199e1933e9db3cb
% canonical-lessons: UR-C212-istri-karna, UR-C212-pisna, UR-C212-talna, UR-C212-ubalna, UR-C212-chakhna

\chapter{To iron, To grind, To fry, To boil, To taste}
\label{ch:ur-to-iron-to-grind-to-fry-to-boil-to-taste}
\hlchaptermodality{212}

\begin{chapteropening}
\textbf{By the end of this chapter:} \emph{I can say to iron, to grind, to fry, to boil and to taste.}

Complete the last lesson of chapter 212: I can say to iron, to grind, to fry, to boil and to taste.
\end{chapteropening}

\section[\texorpdfstring{\ur{استری} \ur{کرنا} (istrī karnā)}{استری کرنا (istrī karnā)}]{\ur{استری} \ur{کرنا} (istrī karnā) — to iron}

\label{lesson:UR-C212-istri-karna}



\begin{quote}
Before the new one: say the Urdu for a servant, then the Urdu for a passer-by.
\end{quote}

\subsection*{You'll want to know: \ur{استری} \ur{کرنا}}
\textbf{\ur{استری} \ur{کرنا}} — \emph{istrī karnā} — \textquotedblleft{}to iron\textquotedblright{}.

\begin{cousinweb}
\textbf{\ur{استری} \ur{کرنا}} (\emph{istrī karnā}) joins \emph{istrī}, \textquotedblleft{}an iron\textquotedblright{}, usually traced to Portuguese \emph{estirar}, \textquotedblleft{}to stretch out\textquotedblright{}, and \emph{karnā}, \textquotedblleft{}to do\textquotedblright{}, from Sanskrit \emph{karoti}, \textquotedblleft{}does\textquotedblright{} (CDIAL 2814).
\end{cousinweb}

\begin{grammarlens}[title={What you've built}]
One more word for this chapter.
\end{grammarlens}

\subsection*{Your turn}
\begin{itemize}
\raggedright
  \item \emph{Say it:} \emph{istrī karnā}
  \item \emph{Say it:} \emph{istrī karnā}, once more
  \item \emph{Say it:} say \emph{istrī karnā}
  \item \emph{Recall:} say the Urdu for a servant, then the Urdu for a passer-by, then say \emph{istrī karnā} again
\end{itemize}

\begin{tcolorbox}[breakable,colback=teal!4,colframe=teal!35!black,title={Before you move on}]
What does \ur{استری} \ur{کرنا} mean? (\textquotedblleft{}To iron\textquotedblright{}.) Say it once more.
\end{tcolorbox}
\section[\texorpdfstring{\ur{پیسنا} (pīsnā)}{پیسنا (pīsnā)}]{\ur{پیسنا} (pīsnā) — to grind}

\label{lesson:UR-C212-pisna}



\begin{quote}
Before the new one: say the Urdu for a passer-by, then the Urdu for to iron.
\end{quote}

\subsection*{You'll want to know: \ur{پیسنا}}
\textbf{\ur{پیسنا}} — \emph{pīsnā} — \textquotedblleft{}to grind\textquotedblright{}.

\begin{cousinweb}
\textbf{\ur{پیسنا}} (\emph{pīsnā}) is a tadbhava, a word inherited from Sanskrit and reshaped by sound change. It comes from \emph{piṃṣati}, \textquotedblleft{}grinds\textquotedblright{} (CDIAL 8142).
\end{cousinweb}

\begin{grammarlens}[title={What you've built}]
One more word for this chapter.
\end{grammarlens}

\subsection*{Your turn}
\begin{itemize}
\raggedright
  \item \emph{Say it:} \emph{pīsnā}
  \item \emph{Say it:} \emph{pīsnā}, once more
  \item \emph{Say it:} say \emph{pīsnā}
  \item \emph{Recall:} say the Urdu for a passer-by, then the Urdu for to iron, then say \emph{pīsnā} again
\end{itemize}

\begin{tcolorbox}[breakable,colback=teal!4,colframe=teal!35!black,title={Before you move on}]
What does \ur{پیسنا} mean? (\textquotedblleft{}To grind\textquotedblright{}.) Say it once more.
\end{tcolorbox}
\section[\texorpdfstring{\ur{تلنا} (talnā)}{تلنا (talnā)}]{\ur{تلنا} (talnā) — to fry}

\label{lesson:UR-C212-talna}



\begin{quote}
Before the new one: say the Urdu for to iron, then the Urdu for to grind.
\end{quote}

\subsection*{You'll want to know: \ur{تلنا}}
\textbf{\ur{تلنا}} — \emph{talnā} — \textquotedblleft{}to fry\textquotedblright{}.

\begin{cousinweb}
\textbf{\ur{تلنا}} (\emph{talnā}) is a tadbhava, a word inherited and reshaped by sound change. It comes from \emph{*talati}, \textquotedblleft{}fries\textquotedblright{}, a form CDIAL reconstructs from its descendants (CDIAL 5736).
\end{cousinweb}

\begin{grammarlens}[title={What you've built}]
One more word for this chapter.
\end{grammarlens}

\subsection*{Your turn}
\begin{itemize}
\raggedright
  \item \emph{Say it:} \emph{talnā}
  \item \emph{Say it:} \emph{talnā}, once more
  \item \emph{Say it:} say \emph{talnā}
  \item \emph{Recall:} say the Urdu for to iron, then the Urdu for to grind, then say \emph{talnā} again
\end{itemize}

\begin{tcolorbox}[breakable,colback=teal!4,colframe=teal!35!black,title={Before you move on}]
What does \ur{تلنا} mean? (\textquotedblleft{}To fry\textquotedblright{}.) Say it once more.
\end{tcolorbox}
\section[\texorpdfstring{\ur{ابالنا} (ubālnā)}{ابالنا (ubālnā)}]{\ur{ابالنا} (ubālnā) — to boil}

\label{lesson:UR-C212-ubalna}



\begin{quote}
Before the new one: say the Urdu for to grind, then the Urdu for to fry.
\end{quote}

\subsection*{You'll want to know: \ur{ابالنا}}
\textbf{\ur{ابالنا}} — \emph{ubālnā} — \textquotedblleft{}to boil\textquotedblright{}.

\begin{cousinweb}
\textbf{\ur{ابالنا}} (\emph{ubālnā}) is a tadbhava, a word inherited and reshaped by sound change. It comes from \emph{*ubbāl}, \textquotedblleft{}make boil\textquotedblright{}, a form CDIAL reconstructs from its descendants (CDIAL 2339). Morgenstierne traced it to \emph{ud-}, \textquotedblleft{}up\textquotedblright{}, and an Indo-European root \textquotedblleft{}to boil\textquotedblright{}; CDIAL finds the forms hard to fit and notes other possible sources.
\end{cousinweb}

\begin{grammarlens}[title={What you've built}]
One more word for this chapter.
\end{grammarlens}

\subsection*{Your turn}
\begin{itemize}
\raggedright
  \item \emph{Say it:} \emph{ubālnā}
  \item \emph{Say it:} \emph{ubālnā}, once more
  \item \emph{Say it:} say \emph{ubālnā}
  \item \emph{Recall:} say the Urdu for to grind, then the Urdu for to fry, then say \emph{ubālnā} again
\end{itemize}

\begin{tcolorbox}[breakable,colback=teal!4,colframe=teal!35!black,title={Before you move on}]
What does \ur{ابالنا} mean? (\textquotedblleft{}To boil\textquotedblright{}.) Say it once more.
\end{tcolorbox}
\section[\texorpdfstring{\ur{چکھنا} (chakhnā)}{چکھنا (chakhnā)}]{\ur{چکھنا} (chakhnā) — to taste}

\label{lesson:UR-C212-chakhna}



\begin{quote}
Before the new one: say the Urdu for to fry, then the Urdu for to boil.
\end{quote}

\subsection*{You'll want to know: \ur{چکھنا}}
\textbf{\ur{چکھنا}} — \emph{chakhnā} — \textquotedblleft{}to taste\textquotedblright{}.

\begin{cousinweb}
\textbf{\ur{چکھنا}} (\emph{chakhnā}) is a tadbhava, a word inherited and reshaped by sound change. It comes from \emph{*cakṣati}, \textquotedblleft{}tastes\textquotedblright{}, a form CDIAL reconstructs from its descendants (CDIAL 4557).
\end{cousinweb}

\begin{grammarlens}[title={What you've built}]
One more word for this chapter.
\end{grammarlens}

\subsection*{Your turn}
\begin{itemize}
\raggedright
  \item \emph{Say it:} \emph{chakhnā}
  \item \emph{Say it:} \emph{chakhnā}, once more
  \item \emph{Say it:} say \emph{chakhnā}
  \item \emph{Recall:} say the Urdu for to fry, then the Urdu for to boil, then say \emph{chakhnā} again
\end{itemize}

\begin{tcolorbox}[breakable,colback=teal!4,colframe=teal!35!black,title={Before you move on}]
What does \ur{چکھنا} mean? (\textquotedblleft{}To taste\textquotedblright{}.) Say it once more.
\end{tcolorbox}
